\begin{frontmatter}

\title{Procurement Rules as Cartel Facilitators: Evidence from Brazil's Minimum-Bidder Requirement\texorpdfstring{\footnote{This version: April 2026. We thank [acknowledgments]. The usual disclaimer applies.}}{}}

\author[insper]{Darcio Genicolo-Martins\corref{cor1}}
\ead{darciogm1@insper.edu.br}
\author[insper]{Paulo Furquim de Azevedo}
\affiliation[insper]{organization={INSPER},
  city={S\~ao Paulo}, country={Brazil}}
\cortext[cor1]{Corresponding author.}

\begin{abstract}
Procurement rules designed to ensure competition can facilitate the
very collusion they aim to prevent. We study Brazil's minimum-bidder
requirement, which mandates at least three participants in sealed-bid
(\emph{convite}) tenders. Using 4.5 million tender-items from S\~ao
Paulo's electronic procurement platform (2009--2019), we show that
this rule creates mechanical demand for cover bidders---firms that
submit deliberately non-competitive bids to satisfy the participation
threshold. Among tenders flagged by frequent-loser presence, 39.2\%
have fewer than three genuine competitors, forcing cartel deployment
of cover firms; in 23.5\%, every bidder is a cover firm, producing
zero genuine competition at positive administrative cost. The price
effect of cover-bidder presence is 3.8\% in rule-bound convite
tenders but 9.3\% in preg\~ao (electronic auction) tenders where the
rule does not apply---a pattern consistent with the rule diluting
rather than generating the collusive signal. An enforcement-capacity
gradient shows the price effect varies 12.5-fold across purchasing
units, concentrating where oversight is weakest. Back-of-the-envelope
welfare calculations suggest the rule costs R\$74--211 million over
the sample period. The findings illustrate a general tension in
procurement design: participation floors intended to protect
competition can instead subsidize the organizational costs of
bid-rigging cartels.
\end{abstract}

\end{frontmatter}

\noindent\textbf{Keywords:} public procurement, bid rigging, cover bidding, minimum-bidder rules, institutional design, cartels

\noindent\textbf{JEL No.:} D44, D73, H57, K21, L41

\bigskip
